\documentclass[../../main-detail.tex]{subfiles}
\begin{document}
\ifSubfilesClassLoaded{\chapter{音韻と造語：制作}}{}

\section{制作のモチベーションとアプローチ}
造語とは，概念に音素列を割り当てる制作である．コノメノでは，形態素未満の音列も含め，音の印象に基づく造語を中心に置く．音韻体系はその素材と制約を与えるが，現行の音素目録や規則のすべてが以下の目標から導かれたわけではない．制作中に見つかった不満は，語の選択だけでなく体系の改訂にも戻す．

\subsection{語彙全体の表情}
語彙全体に固有の表情があり，日常的に使ってもその表情が保たれることを目指す．語の良さには，音列そのものの語感，意味との適合，語彙内での位置，文中での使いやすさが関わる．概念の特徴をよく表す語でも，頻用するには長すぎたり，同じ分野の語を似た顔にしてしまったりする．一語ごとの説明の巧さだけでは評価できない．

\paragraph{検討中の設計目標} 個人的な音の印象を，繰り返しの使用に耐え，日常から専門的思考までを支える語彙へ育てるという方向を前向きに検討する（2026-09-07，定着は未確認）．具体的には，次を制作上の要求とする案である．
\begin{itemize}
    \item 発案時の連想から離れても，繰り返し発音し，手元に置いておきたい形にする．
    \item 日用品や感覚的な対象に加え，退屈な報告，言い淀み，専門的な思考にも使える語彙の広さを持たせる．
    \item 使用のたびに作者の連想の鑑賞を要求せず，普通に使ったり，別の連想を重ねたりできる余地を残す．
\end{itemize}
この目標は現在の感覚をそのまま保存するためのものではない．自然に出る発音は出発点になるが，それだけを完成目標とせず，その響きで会話や思考を続けたいかを問う．

\paragraph{感覚的な接し方} 語形が，対象を別の仕方で感じる入口になることも検討する．典型的な特徴を重ねるだけでなく，対象を捉える視点を選んで音にする．ただし，作れるのはそうした接し方を誘う候補までであり，語を覚えることで実際に見え方が変わるかは使用を通して確かめる必要がある．連想の内容をそのまま語義へ加えるわけではない．

\subsection{統語・文字との連携}
\textbf{[検討中]} 文のどの部分を目立たせ，どの部分を接続や境界として支えるかを，文法上の働き，音のまとまり，字形の構図の間で共同設計する．同じ美意識を掲げるだけでなく，一方の選択が他方の制作条件を変える関係を目指す．意味論の形式体系から特定の音素を導くという主張ではない．

例えば，機能語の背景化には母音や子音の選択だけでなく，独立した韻律的なまとまりを作るかどうかが関わる．文字側では，有声子音を地面の三重線で表すため，音で背景化を狙った有声化が字形の線を増やすこともある．両媒体の効果が一致するとは限らない．具体案と現行体系との衝突は「観察と未解決点」に置く．

\section{語形の印象と構成}\label{sec:zouko-material}
音素や音列の印象を素材とし，リズム，色，語彙の層，分布などから語形の構成を考える．局所的に使える素材と，共通する説明を併存させる．少数の法則への還元を完成条件にはしない．

\memo{まとまりがない．言いたいことがあまり整理されていない．}

\subsection{音素の印象素性}\label{subsec:phoneme-impression}
\noindent
以下は作者個人の印象を整理したものである．造語のフィードバックを受けて調節する．
\memo{重さと物質相はある程度同じ指標か？}
\setlength\LTleft{\fill}%
\setlength\LTright{\fill}%
\begin{longtable}{ccccccccl}
    \caption{母音の印象素性}\\
    \toprule
    音素 & 大きさ & 重さ & 質感 & 色 & 形 & 向き & 物質相 & 動作(自由記述)\\
    \midrule
    \endfirsthead

    \toprule
    音素 & 大きさ & 重さ & 質感 & 色 & 形 & 向き & 物質相 & 動作(自由記述)\\
    \midrule
    \endhead

    \bottomrule
    \endfoot

    \bottomrule
    \endlastfoot

    \lt{a} & 大 & 軽 & 粗 & 黄 & -- & 外 & 気 & 活発に広がる \\
    \lt{o} & 大 & 重 & 滑 & 白 & -- & -- & 固 & 湿気を持ってまとまる \\
    \lt{e} & 中 & 軽 & 粗 & 青 & -- & -- & & 侵蝕的に流れる \\
    \lt{i} & 小 & 軽 & 滑 & 黒 & -- & (内) & & 速く線的に動く \\
    \lt{y} & 小 & 重 & 滑 & 赤 & -- & -- & & 柔らかく有機的に動く \\
    \lt{u} & 小 & 軽 & -- & 緑 & -- & -- & & 流れに任せて流れる \\
    \lt{v} & 小 & 重 & -- & 緑 & -- & -- & & 固まって動かない \\
\end{longtable}
\begin{longtable}{cccccccccl}
    \caption{子音の印象素性}\\
    \toprule
    音素 & 大きさ & 重さ & 質感 & 色 & 形 & 向き & 物質相 & 硬さ & 動作(自由記述)\\
    \midrule
    \endfirsthead

    \toprule
    音素 & 大きさ & 重さ & 質感 & 色 & 形 & 向き & 物質相 & 硬さ & 動作(自由記述)\\
    \midrule
    \endhead

    \bottomrule
    \endfoot

    \bottomrule
    \endlastfoot

    \lt{f} & 小 & 軽 & 粗 & 橙 & 線・面 & -- & & 柔 & 布で包む \\
    \lt{p} & 小 & 軽 & 粗 & -- & -- & 外 & & -- & 大きく広がる \\
    \lt{b} & 中 & 重 & 粗 & 橙 & 線・面 & -- & & 柔 & 小さく跳ぶ \\
    \lt{m} & 中 & 重 & 滑 & 橙 & & -- & & 柔 & 動かずに熟成する \\
    \lt{w} & 中 & 軽 & 滑 & 青 & 面 & -- & -- & -- & 穏やかに広がる \\
    \lt{s} & 中 & 軽 & 粗 & 黄 & 平面 & & & -- & 横方向に速く広がる \\
    \lt{z} & 小 & 軽 & 粗 & -- & & & & -- & 侵蝕的に広がる \\
    \lt{ts} & 小 & 軽 & 粗 & -- & & & & 硬 & 一点に衝撃を与える \\
    \lt{t} & 小 & 軽 & 粗 & -- & 直線 & & & 硬 & 高く跳ねる \\
    \lt{d} & 大 & 重 & 粗 & -- & & & & 硬 & 目的地まで落下する \\
    \lt{n} & 中 & 軽 & 滑 & -- & & & & 硬 & 沈み込む\\
    \lt{l} & 小 & 軽 & 滑 & 黒 & 線 & -- & & -- & 線的に揺れる \\
    \lt{c} & 小 & 軽 & 粗 & 緑 & & & & 柔 & 集団で広がる \\
    \lt{zc} & 大 & 重 & 粗 & 緑 & & & & 柔 & 集団で広がる\\
    \lt{tc} & 中 & 軽 & 粗 & -- & & & & -- & \\
    \lt{j} & 小 & 重 & 滑 & 赤 & & & & 柔 & ウネウネと動く \\
    \lt{kh} & 大 & 軽 & 粗 & -- & & & 気 & -- & (葛藤する) \\
    \lt{k} & 中 & 軽 & 粗 & -- & & & & 硬 & (自在に動く) \\
    \lt{g} & 大 & 重 & 粗 & -- & & & & 硬 & (自在に動く) \\
    \lt{h} & 中 & 軽 & 粗 & 黄 & -- & -- & 気 & -- & 外から包む \\
    \lt{ng} & 中 & 重 & 滑 & -- & & & & -- & \\
\end{longtable}
\setlength\LTleft{0pt}%
\setlength\LTright{0pt}%

母音の印象については，以下のような傾向が見受けられる．
\begin{enumerate}[label=(\arabic*)]
    \item 大きさは開口度に支配されている．
    広い開口$\to$大きい，小さい開口$\to$小さい．
    \item 重さは円唇性に支配されている．
    円唇$\to$重い，非円唇$\to$軽い．
    \item 質感は明確な素性に対応しない．%
\end{enumerate}

子音の印象については，以下のような傾向が見受けられる．
\begin{enumerate}[label=(\arabic*)]
    \item 質感は共鳴音/阻害音の区別に支配されている．
    共鳴音$\to$滑らか，阻害音$\to$粗い．特に摩擦音は強く粗い印象を与える．
\end{enumerate}

\subsection{音素の相性}\label{subsec:phoneme-affinity}

以下は従来のCVの使用頻度と感性をもとにした相性の整理であり，音韻上の許容表ではない．頻度は既存の制作習慣も反映するため，頻度が高いことだけを良さの根拠にはしない．
\begin{table}[H]
    \centering
    \begin{tabular}{cccccccccc}
        \toprule
        C$\to$V & f & m & s & n & t & l & c & h & k\\
        \midrule
        a & ◎ & △ & ◎ & △ & △ & ◯ & ◎ & ◎ & ◯\\
        o & ◯ & △ & △ & ◎ & ◎ & ◎ & ◯ & △ & ◯\\
        e & ◯ & ◎ & △ & △ & △ & ◯ & ◯ & △ & △\\
        \midrule
        i & ◯ & ◯ & × & △ & × & ◯ & ◯ & × & ◯\\
        u & △ & × & ◯ & × & × & ◯ & ◯ & × & △\\
        \bottomrule
    \end{tabular}
\end{table}
ただし，i, u は使用例が少ないため，相性を実際の頻度より増幅して書いている(a, o, eとは前提が違うかもしれない)．ここから次のように考察される：
\begin{enumerate}
    \item 次のように子音を分類すると良い（実例でしっかり反映されているかは難しい）：
    \begin{itemize}
        \item k, f, (m)：やや寛容で\lt{i}を許容する．\lt{m}は\lt{e}を好み，\lt{f}は\lt{k}と\lt{m}の中間の振る舞いを示す．（\lt{m}については不明確だが）軽快な印象．
        \item l, c：寛容で\lt{u}を許容する．有機的な印象．
        \item s, h：\lt{a}と相性が良い（典型的な摩擦音）．風みたいな印象．
        \item n, t：\lt{o}と相性が良い．堅実な印象．
    \end{itemize}
\end{enumerate}

\subsection{音列と意味範囲}
造語対象が与えられた際の音列の割り当ては，音素ごとの印象や短距離の繋がりに加えてさまざまな要素が関わる．
\begin{enumerate}
    \item 意味範囲の確立：（意味が説明しにくい場合もあるが）形態素としてほとんど確立した単位がある（例：\{tent\}「建物」, \{fon\}「袋」, \{kaf\}「芸術」, \{liita\}「楽器」, \{sov\}「原」, \{noi\}「認識」 など）．こうした単位は，より基本的な単位から考察される意味的な適合を覆すことがある．また，単一音素やCV音節くらいの基本的な単位についても（厳密ではないが）質感・動作以上に特定の意味範囲に用いられるという印象がある（例：\lt{la}「食べ物」, \lt{c}「社会」「人間」, \lt{ae}「精神」, \lt{ka}「木製」など）．これらは，明示的な規則ではなく，辞書の「音列」や実際の単語の観察によって理解されるものだと思われる．
    \item 媒体による差異：声に出して読むとき，ラテン文字綴りで書くとき，マリータスユで書くとき，キリル文字・カタカナ転写で書くときなど，媒体や表記体系によって印象は異なる．評価時にはどの媒体による印象かを記録する．媒体ごとの寄与を区別したうえで，音と字形を共同で設計することも検討する．
    \item 作者の母語や文化圏による印象：意味的な適合は，暗黙に作者の母語や文化圏の影響を受けることがある．
    \begin{itemize}
        \item オノマトペの影響：kolon「ビー玉」(コロンッというオノマトペを想起)，danknis「ハンマー」(ダンッというオノマトペを想起)
        \item 単語の影響：造語対象を意味する・関連する既存の単語と近い（意味的適合に寄与），あるいは造語対象とは近くないが既存の単語に近い（語感に寄与する）．これもより基本的な単位から考察される意味的な適合を覆すことがある．例：maton「物質」(マテリアルを想起), plasko「模型」(プラスチックやフラスコを想起), kaci「色」(菓子を想起)．
    \end{itemize}
\end{enumerate}

\subsection{リズム・色・質感}

\subsection{語彙の層と分布}
従来の語彙について，次の傾向が記されていた．これは未集計の観察であり，今後の語を拘束する規則ではない．動物語についても\ko{glisbavn}や\ko{zcagnot}は閉音節で終わるため，一律には当てはまらない．
\begin{enumerate}
    \item 基礎語彙：CVの2--3回の繰り返しが多い．
    \item 動物：開音節で終わる．\lt{$C$i}で終わることが他に比べて多い．
    \item 動作：二重母音（特に\lt{$x$i}，\lt{$x$y}，\lt{$x$v}）あるいは\lt{ja}，\lt{je}で終わる．
    \item 学術語彙：閉音節（特に阻害音のコーダ）で終わる．無声音の長い子音クラスタを用いる．典型的な接続を用いる．
\end{enumerate}

日常語と専門語では使用頻度や周囲の語が異なるため，語彙全体での配置も評価する．\ko{glisbavn}\sentence{牛}は印象に合う一方，使用時には長さが気になり，\ko{zcagnot}\sentence{象}などと大型哺乳類のテイストが似通う．意味的適合を強めると語彙全体がよくなる，とは言えない．

音素・音節・連接・語長などの分布は，制作習慣の偏りや未使用のパターンを探す資料にする．辞書項目内の頻度と発話中の頻度を区別し，参照する自然言語の分布に近づいたことだけを審美的な改善としない．\kotr{oiban}{卵}の語頭\lt{o}，\kotr{stcilkant}{形式的対象}の語頭\lt{stc}は，作者が未発掘のパターンを開拓した例である．

\section{制作と改訂の方法}
意味と音列のどちらを先に決めるかで，制作の方法が変わる．
\begin{enumerate*}
    \item \emph{意味先行}：意味を先に決め音列を後から作る．
    \item \emph{音列先行}：音列を先に作り意味を後から当てる．
\end{enumerate*}
難易度は音列先行の方が低く，特に要請がなければ音列先行を採用するのが良い．意味先行は，理論や実践からの造語要請や既存語の修正，関連語の作成などで必要になる．

\subsection{経路}
\textbf{[暫定：2026-09-15]} 候補を作る経路を次の4つに分ける．2026-09-11から15の制作（家電・身体・材料・日用品・次元類・上位概念，計約120候補）への作者判定から整理した．
\begin{enumerate}
    \item \emph{印象造語}：具体的な場面（色・質感・運動・音）を1つ選び，音素の印象と相性から語根を作る．典型属性の羅列や音素ごとの逐語説明はしない．
    \item \emph{借用改造}：参考言語（モンゴル語，ロシア語，テュルク諸語，ウラル諸語など）の語を，綴りではなく発音から改造する．意味移動を許す（\sentence{団扇}$\to$\kotr{opaakhal}{扇風機}）．
    \item \emph{形態素合成}：既存の語根・接辞を組み合わせる．量接尾辞\lt{-it}（\ko{stoit}，\ko{skolit}，\ko{kulmit}），生起物の\lt{-stav}（\ko{konstav}，\ko{saistav}），ハイフン接頭辞\lt{kon-}／\lt{calf-}（\ko{kon-fenant}）が既に事実上の形態素として機能している．
    \item \emph{音列先行}：借用または確率過程で音列を先に作り，意味を後から当てる．
\end{enumerate}
辞書の「音列」は形態素の辞書を兼ねる（「音列と意味範囲」節）．意味範囲の確立した音列（\{fon\}，\{lek\}「構造」など）は形態素合成に使う．確立していない音列を並べるだけで語を作ることは経路にしない．結果が\sentence{音機}のような複合の読みに落ちるからである（2026-09-30）．偶然の重なりは許す．

\subsection{優先順位}
\textbf{[暫定：2026-09-15]} どの経路を優先するかは，対象の語彙の層で決める．
\begin{itemize}
    \item \emph{日常語}（家電・身体・生物・日用品・材料）：印象造語を主にする．印象は語頭の1音節に置き，2音節目に情報を担わせない．評価の良かった語は単音節（\ko{tcilk}，\ko{pliisk}）か，F音節が重く2音節目が短い語（\ko{weelsa}，\ko{jommat}）であり，2音節目で説明した語は「どれでもよい」と判定された．
    \item \emph{分類語・次元類}（上位概念，量の類）：形態素合成を主にする．日本語でも複合的な概念を単一語根に押し込むと，長く説得力のない語になる．次元類は1音節の程度語根＋\lt{-it}（\ko{pliisk}$\to$\ko{pliiskit}\sentence{輝度量}），生起物の下位類は修飾語根＋\lt{-stav}，次元の数による分割は\lt{kon-}／\lt{calf-}を当てる．
    \item \emph{借用改造}は音列先行に向く．意味先行で使うと，原語の形が対象の軸（大きさ・運動の向き・内部空間・相）と食い違いやすく，印象造語より採択率が低かった（第2セット：印象5/12，借用2/9）．意味先行では，軸に合う形が見つかったときに1案だけ加える．
    \item 候補は「改造の起点」として評価する．採択された語の多くに作者の改造が入る（\ko{opakhal}$\to$\ko{opaakhal}，\ko{tsilk}$\to$\ko{tcilk}）．完成語だけを成果と数えない．
\end{itemize}

\subsection{手順（意味先行）}
\begin{enumerate}
    \item \emph{語義と配置}：語義を1--2文で固定し，上位語と隣接語を辞書IDで挙げる．語義が未決の概念は造語しない．
    \item \emph{軸プロフィール}：人との大きさ比，重さ，質感，運動の向き（内向き／外向き），内部空間の有無，相（固体・液体・気体），丸さを1行で書く．判定で指摘されたズレはこの軸で記録する．
    \item \emph{候補}：概念ごとに3個，経路を2種以上混ぜる．綴りは最初から長ライムにし，F型既定位置の音節が短母音＋阻害音コーダ／開音節にならないようにする（発音時の長音化に頼らない）．
    \item \emph{検査}：\texttt{kono-phonology/scripts/check\_candidates.py}（\texttt{konophon}）で結合表と登録制約，綴りの規則（口蓋化）での書き換え，F型の既定位置での長音化の要否，全存命項目との綴り字Levenshtein距離1以下，配置先の部分木内での語頭2音の共有，語頭オンセットの既使用数，「音列」項目（3文字以上）との重なりを報告する．距離1と語頭共有の候補は差し替える．
    \item \emph{出力}：日付ログの議論用ファイルに，語義・配置・候補表（綴り，音節・型，経路，場面，気になる点，作者判定欄）を書く．体系Pの使用文は付けない．辞書には登録せず，登録は作者の判定後にその指示で行う．
    \item \emph{裁定後}：判定の理由を「観察と未解決点」に戻し，切り出せる部分効果は「音列」項目へ記録する．
\end{enumerate}

\subsection{音列先行}
\begin{enumerate}
    \item 借用ベース
    \item 確率過程ベース
\end{enumerate}
\todo{音列先行の手順を決める．}

\section{観察と未解決点}
\paragraph{機能語への有声音の使用} ミニマルセットは無色ではなくむしろ高明度・彩度を持つ．有声音は従来の大きい・重い・粗いという目立つ印象よりも，色が暗く無個性である印象で捉え直せるのではないか．特に機能語への有声音の使用が自然に思える．
\paragraph{口蓋子音＋\lt{i}語幹と子音接辞の分断} \ko{komatci}に\ko{+kof}を直接付けた\ko{komatcikof}は，\ko{komatc}と\ko{kof}の2単語のように聞こえる．しかし，\ko{kof}は単独では成立しないため，\ko{komatcikof}に違和感がある．修復には次の案がある．
\begin{enumerate}
    \item 核の質を変える：\lt{i}を\lt{y}に変える（例：\ko{+ykof}にする）．
    \item 核の量を増やす：\lt{i}を\lt{ii}にする．
    \item 境界を正当化する：接辞を語として読める長さにし（例：\texttt{nokof}），複合語として分断を許容する．
\end{enumerate}

\paragraph{色彩の平衡化に対する懐疑} 母音調和に似た仕組みである色彩の平衡化は，個別の違和感に対して良い説明を与えるとは限らない．今のところ，平衡化を導入する動機は薄い．平衡化が適用できるかの判断は，その事例で語幹や接辞の母音交替によって自然さが改善するかで行う．

\paragraph{語彙力不足感} \kotr{setayno}{強い}は\kotr{setay}{勝つ}に由来．\kotr{zaalonik}{ラジオ}は「ザー」からの連想に由来．これらは``勝つみたいな''，``ザーっていう機械''（合わせて``$A$的な〜''）語彙力不足な説明がダサい．

\paragraph{場違い感} 語彙の印象に過度に合わせた結果，辞書全体の中で場違いに聞こえる場合がある．\kotr{boni}{醜い}，\kotr{gulu}{死体}は意味に合わせて音も汚く作られているが合わない．\kotr{lizlim}{チョーカー}はメンヘラ感があるが合わない．

\paragraph{母音融合と混色} 旧発音で\lt{ea}を\textipa{[ea]}と読むと青と黄の主張が強く，新発音で\textipa{[\ae]}としてまとめると灰色寄りに感じられる（青と黄が補色で混色は無彩色になる立場）．

\paragraph{カテゴリを表す形態素の音節構造} CV，CVC，CVVはダサく感じる傾向がある（例：\{na\}，\{nis\}）．CCVCやCVCCのように子音が重い方が良い（例：\{tent\}（施設），\{stakk\}（数））．ただし，\{fon\}のように平気なのもある．

\end{document}
